\chapter{Conclusion}
\label{chap:conclusion}

\section{Summary of Achievements}
\label{sec:summary_achievements}

This thesis investigated, engineered, and empirically validated an end-to-end passive monocular RGB face anti-spoofing framework designed to protect CPU-based biometric attendance verification against two-dimensional presentation attacks. Motivated by the vulnerability of standard angular-margin face recognition models to smartphone video replays and printed photographs, every stage of the system---from person-disjoint dataset curation and normalization diagnostics to frozen-graph CPU compilation and WebSocket attendance gating---was developed and audited against the objectives established in Chapter~\ref{chap:introduction}.

First, at the data and benchmark level, progressive audits exposed how random frame-level splits, unmasked presenter faces in wide-angle recordings, and cross-container detector mismatches inflate laboratory accuracy while failing on unseen subjects. To resolve these confounders, an $89$-video Merged Multi-Source Benchmark ($28{,}957$ crops) and a dedicated $91$-video target-sensor benchmark (\texttt{Cam4K}, $26{,}094$ crops extracted from $140{,}707$ raw SC-M200 frames) were constructed under strict person-disjoint partitioning, six-stage optical quality filtering, and automated right-edge presenter masking.

Second, at the representation and optimization level, controlled ablations demonstrated that combining imbalanced video cohorts with scalar loss weights ($52.75\times$ ratio) corrupts unweighted Batch Normalization moving statistics, doubling spoof false acceptance ($\text{APCER}$) from $7.88\%$ to $16.16\%$. Introducing persistent $4\times8$ $(source, class)$ stratified mini-batch queues alongside sample-local Group and Layer Normalization reduced merged-test $\text{ACER}$ to $4.82\%$ and eliminated false session verdicts. On the target SC-M200 sensor, adapting the $545{,}154$-parameter \texttt{LivenessNet} classifier head atop a frozen universal convolutional backbone with WiSE-FT weight-space interpolation achieved $98.09\%$ frame accuracy, $\text{APCER} = 3.76\%$, $\text{BPCER} = 0.12\%$, and $\text{ACER} = 1.94\%$ on unseen test subjects, while classifying $100.00\%$ ($14/14$) of test videos correctly and retaining $99.84\%$ accuracy across $600$ external laptop-webcam crops.

Third, at the systems and deployment level, compiling \texttt{LivenessNet} into an unquantized LiteRT flatbuffer executed via XNNPACK reduced single-crop CPU latency to a median of $0.163$~ms ($198\times$ faster than eager Keras execution). Coupled with YuNet ONNX face detection, greedy $\text{IoU}$ multi-face tracking, recency-weighted \texttt{FastBinaryDecision} stabilization, and $6$-observation fail-closed \texttt{TemporalDecision} hysteresis, the shared runtime sustained $30.4$~FPS in multi-threaded live camera monitoring and $40.61$~ms median end-to-end latency in Waitress HTTP serving. Finally, integrating the YuNet-realigned \texttt{LivenessGate} upstream of calibrated ArcFace prototype matching ($\tau = 0.55$) inside \texttt{ChamCong\_FaceID} deterministically blocked smartphone replay attacks while preserving low-friction check-in logging for enrolled personnel.

\section{Limitations and Future Work}
\label{sec:limitations_future_work}

Despite fulfilling its accuracy and real-time CPU latency targets, the current system exhibits four practical limitations that define clear boundaries for deployment:
\begin{itemize}
    \item \textbf{Presentation attack modality coverage:} Empirical training in \texttt{Cam4K} concentrated primarily on electronic screen replay instruments (smartphones, tablets, and monitors). Although auxiliary $YC_rC_b$/HSV optical descriptors mitigate diffuse paper veils, high-grade dye-sublimation matte prints, three-dimensional flexible silicone masks, and kernel-level virtual camera stream injection were not quantitatively benchmarked.
    \item \textbf{Sample size at the video unit of analysis:} While frame-level evaluation spans $4{,}811$ unseen test crops ($9{,}826$ combined validation and test crops), the independent subject-disjoint test partition comprises $14$ video sequences ($16$ sequences in Cohort~IIa). Larger multi-hundred-subject cohorts across diverse age groups and skin tones are needed to tighten video-level confidence intervals below $1\%$.
    \item \textbf{Environmental and pose envelope:} Operating purely in the visible RGB spectrum leaves passive classification sensitive to extreme low-light conditions ($\mu_{\text{gray}} < 35$), direct sunlight washout, and steep off-axis head rotations exceeding $\pm 35^\circ$ where facial skin area shrinks within the $32\times32$ patch.
    \item \textbf{Purely passive verification ceiling:} High-pixel-density OLED displays positioned carefully to match ambient wall illumination (as observed in \texttt{fake\_REC-0041-A}) narrow single-frame decision margins when subpixel aliasing falls below camera sensor resolution.
\end{itemize}

Overall, these constraints motivate four concrete directions for future research and engineering extension:
\begin{enumerate}
    \item \textbf{Conditional active-passive hybrid verification:} Retaining instant passive verification for clear high-confidence presentations ($\bar{P}_{\text{real}} \ge 0.85$) while triggering a lightweight randomized gaze, blink, or screen-flash reflection challenge only when session confidence falls into an ambiguous band ($0.35 < \bar{P}_{\text{real}} < 0.85$) would strengthen defense against high-PPI replays and 3D masks without slowing routine check-ins.
    \item \textbf{Expanded physical print and 3D mask corpora with frequency-domain branches:} Augmenting \texttt{LivenessNet} with a lightweight Fast Fourier Transform (FFT) or central-difference gradient branch trained on laser, inkjet, glossy photo-paper, and rigid mask presentations can sharpen discrimination on non-emissive substrates.
    \item \textbf{Browser-side WebAssembly / WebGPU inference:} Deploying the compact $2.18$~MB \texttt{liveness\_cam4k.tflite} flatbuffer directly inside client browsers via MediaPipe/LiteRT-Web would offload frame streaming bandwidth from the central attendance server and preserve employee video privacy on local devices.
    \item \textbf{Continual domain adaptation across enterprise cameras:} Incorporating unsupervised test-time batch/instance calibration or periodic federated WiSE-FT anchor updates across newly installed office cameras can maintain long-term calibration without full dataset recollection.
\end{enumerate}
